\chapter{Ileostomy}

\section*{Introduction \& Clinical Indications}
An ileostomy is a surgical procedure that involves creating an opening, called a stoma, on the abdominal surface through which a portion of the small intestine, specifically the ileum, is brought to the exterior. This stoma serves as a diversion for fecal matter, allowing it to bypass the large intestine (colon) and rectum entirely. The procedure is essential for managing various gastrointestinal disorders, including inflammatory bowel disease, colorectal cancer, and traumatic injuries to the bowel. By providing a means for continuous drainage of liquid to semi-liquid stool into an external collection appliance, ileostomies can significantly improve the quality of life for patients facing severe bowel dysfunction. Depending on the clinical scenario, an ileostomy may be temporary, allowing for healing of the distal bowel, or permanent, serving as the definitive route for fecal elimination following total proctocolectomy. The clinical significance of this surgery lies in its ability to alleviate symptoms, prevent complications, and enhance the overall well-being of patients with debilitating gastrointestinal conditions.

\begin{itemize}
  \item Ileal stoma
  \item Bowel diversion surgery
  \item Brooke ileostomy (specific technique)
  \item Loop ileostomy
  \item End ileostomy
  \item Permanent ileostomy
  \item Temporary ileostomy
\end{itemize}

The primary surgical principle of an ileostomy is to divert the fecal stream from a diseased or compromised segment of the large bowel, allowing for healing or providing a permanent alternative route for waste elimination. This is accomplished by transecting the ileum, typically at its terminal portion, and bringing the proximal end through a surgically created opening in the abdominal wall. The exposed segment of ileum is everted and sutured to the skin, forming a stoma that directs effluent away from the skin surface and into an external collection appliance.

The rationale for this diversion is multifaceted. In cases of inflammatory bowel disease or colorectal cancer requiring total proctocolectomy, the ileostomy provides a definitive means of fecal elimination. For conditions such as low rectal anastomoses, severe trauma, or perianal sepsis, a temporary diverting ileostomy protects a distal healing site from fecal contamination, thereby reducing the risk of anastomotic leak and promoting optimal healing. The continuous flow of liquid small-bowel contents into an external appliance prevents accumulation within the bypassed or resected colon and rectum, alleviating symptoms and facilitating recovery. Depending on the underlying condition, the distal bowel segment may be oversewn and left in the abdomen (Hartmann's pouch) or brought out as a mucous fistula.

The practice of diverting the fecal stream has historical roots dating back to ancient medicine, where early attempts to manage bowel obstruction were often rudimentary and fraught with complications. The modern ileostomy began to take shape in the early 20th century, with initial procedures often resulting in significant morbidity due to skin excoriation from the highly enzymatic effluent and the absence of effective collection devices.

A significant advancement in ileostomy technique was made by Dr. Bryan Brooke in the 1950s, who introduced the eversion technique, known as the Brooke ileostomy. This method involved everting the bowel mucosa over the serosa and suturing it to the skin, creating a spouted stoma that significantly reduced peristomal skin complications and improved appliance adherence. This technique became the gold standard for ileostomy creation. Further refinements were introduced by surgeons such as Dr. Rupert Turnbull, who emphasized the importance of meticulous preoperative stoma site marking by specialized ostomy nurses, recognizing that proper placement is crucial for long-term patient success. The development of modern ostomy appliances in the latter half of the 20th century, characterized by improved adhesion, odor control, and discretion, transformed the ileostomy from a last-resort procedure into a manageable and often life-saving intervention, allowing patients to lead active and fulfilling lives.

An ileostomy is indicated for a variety of gastrointestinal conditions, particularly when the large intestine or rectum is diseased, requires removal, or needs to be bypassed for healing.

\begin{itemize}
  \item To treat severe, medically refractory ulcerative colitis requiring total proctocolectomy.
  \item To manage complicated Crohn's disease affecting the colon or rectum, especially with strictures, fistulas, or toxic megacolon, where medical therapy has failed.
  \item As part of a total proctocolectomy for familial adenomatous polyposis (FAP) to prevent colorectal cancer, which is inevitable without surgical intervention.
  \item To divert the fecal stream and protect a distal colorectal or coloanal anastomosis after low anterior resection for rectal cancer, particularly in high-risk anastomoses.
  \item In cases of severe colorectal trauma or perforation where primary repair is not safe or feasible, necessitating immediate fecal diversion.
  \item To decompress an obstructed colon or rectum that cannot be managed by other means, such as in cases of advanced colorectal cancer or severe diverticular disease.
  \item For intractable perianal sepsis, complex rectovaginal fistulas, or severe perianal Crohn's disease that requires complete fecal diversion to facilitate healing.
  \item In rare instances of severe, intractable constipation or colonic inertia unresponsive to maximal medical therapy and other surgical options.
  \item As a temporary measure to allow healing of an anastomotic leak or severe peritonitis following colorectal surgery.
\end{itemize}

An ileostomy can serve both as a primary and a secondary procedure, depending on the specific clinical context. It is considered a primary procedure when it is the definitive surgical intervention for a condition, such as a total proctocolectomy for medically refractory ulcerative colitis or familial adenomatous polyposis. In these scenarios, the entire colon and rectum are removed, and the ileostomy becomes the permanent means of fecal elimination.

Conversely, an ileostomy often functions as a secondary or adjunct procedure. For example, a temporary diverting loop ileostomy is frequently created after a low anterior resection for rectal cancer to protect the newly formed anastomosis from fecal contamination during its initial healing phase. In this context, the ileostomy is a planned, temporary measure, with the explicit intention of subsequent reversal once the anastomosis has healed. It can also be a salvage procedure in cases of anastomotic leak, severe peritonitis, or uncontrolled sepsis following other abdominal surgeries, where immediate fecal diversion is necessary to control infection and stabilize the patient. The decision for primary versus secondary positioning is guided by the urgency of the condition, the extent of disease, and the potential for future bowel reconstruction.

\section*{Relevant Surgical Anatomy}
The ileum is the terminal segment of the small intestine, extending from the jejunum to the ileocecal valve, where it connects with the large intestine. Its primary functions include the absorption of nutrients, vitamin B12, and bile salts. The arterial supply to the ileum is provided by the ileal branches of the superior mesenteric artery (SMA), which form a network of arcades within the mesentery, ensuring adequate blood flow. Venous drainage follows the arterial supply, with blood draining into the superior mesenteric vein (SMV) and subsequently into the hepatic portal system. Lymphatic drainage occurs through mesenteric vessels to the superior mesenteric lymph nodes.

When creating the stoma, the anatomy of the abdominal wall is critical. The stoma is ideally positioned through the rectus abdominis muscle in the right lower quadrant, which provides a supportive muscular framework around the stoma, reducing the risk of parastomal hernia. The layers traversed during stoma creation include the skin, subcutaneous fat, anterior rectus sheath, rectus abdominis muscle, posterior rectus sheath, and peritoneum. Care must be taken to avoid injury to the neurovascular bundles of the abdominal wall, particularly the intercostal nerves (T7-T12) and their accompanying vessels, to minimize postoperative pain and sensory deficits. The chosen stoma site should be free from bony prominences, skin folds, scars, and the beltline to ensure optimal appliance adherence and patient comfort.

\section*{Preoperative Workup \& Preparation}
Thorough preoperative preparation is essential for a successful ileostomy and optimal patient recovery, minimizing complications and facilitating adaptation. This involves a comprehensive medical evaluation, patient education, and specific surgical preparations.

Key aspects of preoperative workup and preparation include:

\begin{itemize}
  \item \textbf{Stoma Site Marking:} This is a critical step, performed by an experienced ostomy nurse. The ideal site is chosen while the patient is sitting, standing, and lying down, considering body contours, skin folds, scars, and clothing lines to ensure optimal appliance fit, patient comfort, and ease of self-care.
  
  \item \textbf{Medical Optimization:} A thorough medical evaluation, including cardiac and pulmonary clearance, is performed to assess the patient's fitness for general anesthesia and major abdominal surgery. Any comorbidities such as diabetes, hypertension, or renal insufficiency are optimized.
  
  \item \textbf{Laboratory and Imaging Studies:} Routine blood tests (complete blood count, electrolytes, renal and liver function, coagulation profile, blood type and cross-match) are performed. Imaging studies (e.g., CT scan, MRI, colonoscopy) are reviewed to confirm diagnosis, assess disease extent, and plan the surgical approach.
  
  \item \textbf{Bowel Preparation:} Depending on the surgeon's preference and the urgency of the surgery, a mechanical bowel preparation (e.g., polyethylene glycol solution) combined with oral antibiotics (e.g., neomycin and metronidazole) may be prescribed to cleanse the colon and reduce bacterial load, particularly if a distal anastomosis is planned.
  
  \item \textbf{Dietary Modifications and NPO Status:} Patients are typically advised to follow a low-residue diet for several days prior to surgery. NPO (nil per os) status is required for at least 6-8 hours before anesthesia to minimize the risk of aspiration.
  
  \item \textbf{Medication Review and Adjustment:} All medications, especially anticoagulants, antiplatelet agents, immunosuppressants, and insulin, must be reviewed and adjusted or held as per surgical and anesthesia guidelines to minimize bleeding risk and manage chronic conditions.
  
  \item \textbf{Prophylactic Antibiotics:} Intravenous broad-spectrum antibiotics (e.g., cefazolin and metronidazole) are administered shortly before the incision to reduce the risk of surgical site infection.
  
  \item \textbf{Informed Consent and Patient Education:} A detailed discussion of the procedure, its risks, benefits, alternatives, and the implications of living with a stoma is conducted. Patients receive extensive education from an ostomy nurse regarding stoma care, appliance management, dietary adjustments, and potential complications. Psychological support and peer counseling may also be offered to help patients prepare for the lifestyle changes.
\end{itemize}

The decision to create an ileostomy involves careful consideration of potential contraindications and necessary precautions to ensure patient safety and optimize outcomes.

\textbf{Absolute Contraindications:}
\begin{itemize}
  \item \textbf{Uncorrectable Severe Coagulopathy or Bleeding Diathesis:} Patients with severe, unmanageable bleeding disorders are at extremely high risk of intraoperative and postoperative hemorrhage.
  
  \item \textbf{Uncontrolled Sepsis or Hemodynamic Instability:} Patients in septic shock or with severe hemodynamic instability are too unstable for major abdominal surgery, and attempts to proceed would carry an unacceptably high mortality risk.
  
  \item \textbf{Severe, Uncorrectable Malnutrition or Cachexia:} Extreme malnutrition significantly impairs wound healing, immune function, and overall surgical tolerance, leading to very high morbidity and mortality rates.
  
  \item \textbf{Extensive Intra-abdominal Adhesions or Diffuse Carcinomatosis:} In some cases, the abdominal cavity may be so obliterated by adhesions from previous surgeries or diffuse cancer that safe mobilization of the bowel and creation of a stoma without significant injury is impossible.
\end{itemize}

\textbf{Relative Contraindications and Precautions:}
\begin{itemize}
  \item \textbf{Active Crohn's Disease in Proposed Ileal Segment:} While often an indication for ileostomy, active Crohn's disease in the specific segment of ileum chosen for the stoma can lead to recurrence, stricture, fistula formation, or poor healing at the stoma site. Careful selection of a healthy, uninvolved segment is crucial.
  
  \item \textbf{Prior Radiation Injury:} Patients who have received prior pelvic or abdominal radiation may have fibrotic, poorly vascularized bowel and abdominal wall tissues, increasing the risk of ischemia, poor healing, and complications such as fistula formation or wound dehiscence.
  
  \item \textbf{Significant Obesity:} Morbid obesity can make stoma creation technically challenging due to thick abdominal wall fat, increase the risk of parastomal hernia, and complicate appliance fitting and adherence, leading to peristomal skin issues.
  
  \item \textbf{Dense Adhesions from Previous Surgeries:} Extensive prior abdominal surgeries can result in dense adhesions, increasing operative time, the risk of iatrogenic bowel injury, and difficulty in mobilizing a tension-free ileal segment.
  
  \item \textbf{Short Bowel Syndrome:} Patients with a history of extensive small bowel resections may be at high risk of severe dehydration, electrolyte imbalances, and malnutrition due to high ileostomy output, necessitating careful fluid and nutritional management.
  
  \item \textbf{Patient Factors (Cognitive Impairment, Lack of Support):} Poor patient compliance, significant cognitive impairment, or lack of adequate social support can hinder effective stoma care, leading to complications and poor quality of life. Extensive preoperative education and support systems are vital.
  
  \item \textbf{Technical Precautions:} Meticulous surgical technique is paramount. This includes avoiding tension on the exteriorized ileum, ensuring adequate blood supply to the stoma (assessed by color and bleeding from the cut edge), creating a properly spouted stoma to prevent skin irritation, and ensuring the trephine is neither too tight nor too wide.
\end{itemize}

\section*{Surgical Technique \& Procedure}
The creation of an ileostomy is a meticulous surgical procedure typically performed under general anesthesia. Patient positioning is supine, and the abdomen is prepped and draped in a sterile fashion. The exact location of the stoma is crucial and is ideally marked preoperatively by an experienced ostomy nurse, usually in the right lower quadrant, within the rectus abdominis muscle, away from bony prominences, skin folds, scars, and the beltline. The procedure can be performed via an open laparotomy or laparoscopically, with the latter offering benefits of smaller incisions and potentially faster recovery.

The key steps involved in ileostomy creation are:

\begin{itemize}
  \item \textbf{Abdominal Access and Bowel Mobilization:} A midline laparotomy incision or several small laparoscopic ports are used to gain access to the abdominal cavity. The small bowel, specifically the terminal ileum, is identified and mobilized. For an end ileostomy, the ileum is divided, and the distal end is either removed with the resected colon or oversewn and left as a Hartmann's pouch. For a loop ileostomy, a segment of ileum (typically 15-20 cm proximal to the ileocecal valve) is selected, and a small incision is made in its mesentery to create a window, allowing for the passage of a plastic rod or bridge later.
  
  \item \textbf{Trephine Creation:} A circular incision (approximately 2.5-3 cm in diameter) is made in the skin at the pre-marked stoma site. The subcutaneous fat is dissected, and the anterior rectus sheath is incised in a cruciate or circular fashion. The rectus muscle fibers are carefully separated longitudinally, and the posterior rectus sheath and peritoneum are incised to create a trephine (opening) through the abdominal wall. This opening must be wide enough to accommodate the bowel without tension or ischemia, typically allowing two fingers to pass comfortably.
  
  \item \textbf{Bowel Exteriorization:} The proximal end of the divided ileum (for an end ileostomy) or the selected loop of ileum (for a loop ileostomy) is gently brought through the trephine. Care is taken to ensure the bowel is not twisted along its mesenteric axis and that its mesentery is not under tension, which could compromise blood supply. The bowel should protrude approximately 2-3 cm above the skin level.
  
  \item \textbf{Stoma Maturation (Brooke Technique):} For an end ileostomy, the ileum is everted, meaning the serosal surface is tucked inward, and the mucosal surface is brought out to form a spout. The everted ileal edge is then sutured to the dermis of the skin incision using absorbable sutures (e.g., 3-0 or 4-0 Vicryl), creating a healthy, spouted stoma that directs effluent away from the skin. For a loop ileostomy, both the proximal (effluent-producing) and distal (mucous fistula) limbs are brought through the trephine, and the loop is matured to the skin, often with a plastic rod or bridge placed temporarily beneath the loop to prevent retraction.
  
  \item \textbf{Abdominal Closure:} The abdominal incision (laparotomy or port sites) is closed in layers. Drains are typically not required for an uncomplicated ileostomy unless there is significant contamination, a large dead space, or a high risk of anastomotic leak in the case of a protective ileostomy. A stoma appliance (pouch) is immediately placed over the newly formed stoma to collect effluent and protect the skin.
\end{itemize}

\section*{Complications \& Risks}
Ileostomy creation, like any major surgical procedure, carries inherent risks. These can be broadly categorized as general surgical complications and specific stoma-related complications.

\textbf{General Surgical Risks:}
\begin{itemize}
  \item \textbf{Bleeding:} Intraoperative or postoperative hemorrhage, which may manifest as hematoma, significant blood loss requiring transfusion, or even reoperation.
  
  \item \textbf{Infection:} Surgical site infection (SSI) at the abdominal incision or around the stoma, intra-abdominal abscess, or generalized sepsis.
  
  \item \textbf{Anesthesia Complications:} Adverse reactions to anesthetic agents, respiratory depression, cardiac events (e.g., myocardial infarction, arrhythmia), cerebrovascular accident (stroke), or deep vein thrombosis (DVT) and pulmonary embolism (PE).
  
  \item \textbf{Bowel Injury:} Accidental injury to other segments of the bowel, bladder, ureters, or major blood vessels during dissection, potentially requiring repair or further surgical intervention.
  
  \item \textbf{Wound Dehiscence:} Partial or complete separation of the abdominal incision, requiring wound care or reoperation.
  
  \item \textbf{Adhesions and Internal Hernia:} Formation of scar tissue within the abdomen, which can lead to future bowel obstruction or internal herniation.
\end{itemize}

\textbf{Procedure-Specific Risks and Complications:}
\begin{itemize}
  \item \textbf{Stoma Ischemia or Necrosis:} Inadequate blood supply to the exteriorized ileum, leading to tissue death. This is a surgical emergency requiring urgent reoperation to revise the stoma.
  
  \item \textbf{Stoma Retraction:} The stoma pulls back below the skin level, making appliance fitting difficult, increasing the risk of peristomal skin excoriation, and potentially leading to intra-abdominal leakage.
  
  \item \textbf{Stoma Stenosis:} Narrowing of the stoma opening, impeding effluent flow and potentially causing symptoms of bowel obstruction, requiring dilation or surgical revision.
  
  \item \textbf{Parastomal Hernia:} Protrusion of abdominal contents (e.g., small bowel, omentum) through the abdominal wall defect adjacent to the stoma. This is a common late complication, often requiring surgical repair.
  
  \item \textbf{Stoma Prolapse:} Protrusion of the bowel through the stoma opening, which can be partial or complete. It can be managed conservatively or surgically, depending on severity and symptoms.
  
  \item \textbf{High-Output Ileostomy:} Excessive fluid and electrolyte loss through the stoma (typically >1000-1500 mL/day), leading to dehydration, severe electrolyte imbalances (especially hyponatremia, hypokalemia, hypomagnesemia), and acute kidney injury.
  
  \item \textbf{Peristomal Skin Excoriation/Irritation:} Damage to the skin around the stoma due to prolonged exposure to digestive enzymes from the effluent, causing pain, itching, redness, and potential infection. This is often due to improper appliance fit or care.
  
  \item \textbf{Fistula Formation:} Abnormal connection between the stoma and another organ or the skin, often a complication of underlying inflammatory bowel disease or technical error.
  
  \item \textbf{Bowel Obstruction:} Can occur early due to adhesions, technical issues (e.g., narrow trephine), or internal herniation, or late due to strictures, adhesions, or recurrent disease.
  
  \item \textbf{Psychological Impact:} Body image issues, anxiety, depression, social isolation, and sexual dysfunction are common psychological challenges that patients may face after ileostomy creation.
\end{itemize}

\section*{Postoperative Care \& Recovery}
The postoperative recovery timeline for an ileostomy patient involves several distinct phases, from immediate post-anesthesia care to long-term adaptation.

Immediately after the procedure, the patient is transferred to the Post-Anesthesia Care Unit (PACU) for close monitoring of vital signs, pain control, and initial assessment of the newly created stoma. The stoma should appear pink or red, moist, and viable. Within the first 24-48 hours, the stoma typically begins to function, producing liquid effluent. Output is carefully measured, and fluid and electrolyte balance is a critical concern due to the potential for high ileostomy output; intravenous fluids are administered to prevent dehydration, and electrolyte levels are monitored daily. Early ambulation is encouraged to prevent complications like deep vein thrombosis and promote bowel function. Pain management is crucial, often involving patient-controlled analgesia (PCA) initially, transitioning to oral pain medications.

Over the first 1-2 weeks, oral intake is gradually advanced from clear liquids to a low-residue diet as tolerated, typically once bowel sounds return and the patient is passing flatus or stool through the stoma. Ostomy nurses play a vital role in patient education, teaching them how to empty and change the appliance, assess stoma health, and manage output. Patients learn about dietary modifications to manage output consistency and minimize gas or odor. Discharge planning focuses on ensuring the patient and their caregivers are proficient in stoma care, understand warning signs of complications, and have access to necessary supplies and support. Full recovery and adaptation to living with an ileostomy can take several weeks to months, involving significant adjustments to diet, lifestyle, and body image. Patients are advised to avoid heavy lifting for 6-8 weeks to reduce the risk of parastomal hernia.

During an ileostomy procedure, the surgeon meticulously documents several key findings. These include the specific segment of ileum used, its vascularity, the tension-free nature of its exteriorization, the size and location of the trephine, and the appearance of the matured stoma. Any intra-abdominal adhesions, the condition of the bypassed or resected colon/rectum, and the presence of any other incidental findings are also noted. For temporary ileostomies, the integrity of the distal anastomosis (if created) is confirmed. The surgeon's findings provide immediate feedback on the technical success of the stoma creation and the overall abdominal exploration.

The pathology report on any resected tissues (e.g., colon, rectum, lymph nodes) is crucial for confirming the underlying diagnosis and guiding further management. For inflammatory bowel disease, the report details the type (e.g., ulcerative colitis, Crohn's disease), severity, and extent of inflammation. For colorectal cancer, it provides the tumor type, grade, depth of invasion, lymph node involvement, presence of vascular or perineural invasion, and margin status. These pathological findings are essential for cancer staging (e.g., TNM classification) and determining the need for adjuvant therapies such as chemotherapy or radiation. The pathology report thus provides the definitive diagnosis and prognostic information, complementing the surgical findings.

A normal and successful postoperative course following an ileostomy procedure is characterized by a predictable progression towards recovery and adaptation. Initially, the patient will experience incisional pain, which should be well-controlled with analgesia. The stoma itself should be pink or beefy red, moist, and viable, indicating good blood supply. Stoma function typically commences within 24-72 hours, with the passage of liquid to semi-liquid effluent into the collection appliance. The volume of output will stabilize over the first few days, usually settling into a range of 500-1000 mL per 24 hours. The peristomal skin should remain healthy, without redness or irritation, signifying proper appliance fit and care. As the patient recovers, they will gradually resume oral intake, ambulate independently, and learn to manage their stoma care with the guidance of ostomy nurses. Resolution of the original symptoms that necessitated the surgery (e.g., abdominal pain, bleeding, obstruction) is expected. Ultimately, a successful course leads to the patient's independence in stoma management, improved quality of life, and the absence of significant early or late complications.

Postoperative care and long-term follow-up are crucial for patients with an ileostomy to ensure proper healing, address any complications, and facilitate adaptation to their new lifestyle.

\begin{itemize}
  \item \textbf{Hospital Discharge and Initial Home Care:} Patients are typically discharged when they are tolerating oral intake, managing pain effectively, and demonstrating proficiency in stoma care. Home health nursing or continued ostomy nurse support is often arranged to provide ongoing education and assistance.
  
  \item \textbf{First Postoperative Visit:} Usually scheduled 1-2 weeks after discharge to assess wound healing, stoma viability, peristomal skin integrity, and address any immediate concerns. Non-absorbable sutures around the stoma or abdominal incision may be removed at this time.
  
  \item \textbf{Ostomy Nurse Consultations:} Ongoing support from an experienced ostomy nurse is vital for troubleshooting appliance issues, providing dietary counseling, managing high output, and offering psychological support. These consultations may continue for several months or longer.
  
  \item \textbf{Hydration and Electrolyte Monitoring:} For patients with persistently high output, periodic blood tests to check electrolyte levels (sodium, potassium, magnesium) and renal function may be necessary to prevent dehydration and electrolyte imbalances.
  
  \item \textbf{Dietary Counseling:} Guidance on a low-residue diet initially, followed by gradual introduction of foods, with specific advice on managing gas, odor, and output consistency. Patients learn to identify foods that may cause issues.
  
  \item \textbf{Activity Resumption:} Gradual return to normal activities is encouraged, but patients are advised to avoid heavy lifting (typically >10-15 lbs) for several weeks to months (e.g., 6-12 weeks) to reduce the risk of parastomal hernia formation.
  
  \item \textbf{Imaging and Reversal Surgery (for temporary ileostomies):} For temporary ileostomies, imaging (e.g., contrast enema or flexible endoscopy) may be performed after several weeks or months (typically 8-12 weeks) to confirm healing and integrity of the distal anastomosis before planning ileostomy reversal surgery.
  
  \item \textbf{Warning Signs:} Patients are educated to contact their healthcare provider immediately for signs of stoma changes (e.g., dark color, severe pain, bleeding), persistent high output with symptoms of dehydration (e.g., dizziness, decreased urination), severe abdominal pain, fever, or signs of bowel obstruction (e.g., nausea, vomiting, no output).
\end{itemize}

Regular follow-up with the surgeon and primary care physician is also important for overall health management and screening.

\section*{Outcomes \& Prognosis}
Ileostomy creation is a relatively common surgical procedure globally, with its incidence closely tied to the prevalence of severe inflammatory bowel disease (IBD) and colorectal cancer. Ulcerative colitis and Crohn's disease are major indications, affecting millions worldwide, with a significant proportion (up to 25-40\% of UC patients and 20-30\% of CD patients) eventually requiring surgical intervention, including ileostomy, over their lifetime. Colorectal cancer, the third most common cancer globally, also frequently necessitates ileostomy, particularly for low rectal cancers where a diverting stoma protects a distal anastomosis or for advanced cases requiring total proctocolectomy. The age distribution for ileostomy creation is typically bimodal, with peaks in young adults (20-40 years) due to IBD and in older adults (60+ years) due to colorectal cancer, diverticular disease complications, or other age-related bowel pathologies. There is no significant sex predilection for ileostomy itself, though the underlying diseases may show slight differences in incidence between sexes. While exact global incidence rates are challenging to ascertain due to varying reporting, it is estimated that hundreds of thousands of ileostomies are performed annually worldwide. Mortality directly attributable to elective ileostomy creation is low (typically <1-2\%), but morbidity, particularly stoma-related complications, can be substantial, affecting up to 50-70\% of patients over their lifetime, necessitating ongoing care and sometimes reoperation.

The outcomes and prognosis following ileostomy creation are generally favorable, especially in elective settings, with significant improvement in quality of life for many patients. For conditions like medically refractory ulcerative colitis, total proctocolectomy with permanent ileostomy is considered curative, eliminating the disease and its associated symptoms, and preventing the risk of colorectal cancer. For colorectal cancer, ileostomy, whether temporary or permanent, plays a crucial role in achieving disease control and improving survival by facilitating definitive resection or protecting a healing anastomosis. Typical success rates for symptom resolution and disease control are high, often exceeding 80-90\%. For example, in patients undergoing ileostomy for ulcerative colitis, the APPAC trial demonstrated the safety and efficacy of laparoscopic ileostomy creation.

However, the long-term prognosis is influenced by the underlying disease, the presence of complications, and the patient's ability to adapt to stoma care. While an ileostomy can significantly improve physical health and survival, some patients experience challenges with body image, social integration, and sexual function, which can impact overall quality of life. Stoma-related complications, such as parastomal hernia (incidence up to 50\% over 10 years), prolapse, or high output, can necessitate further surgical intervention or ongoing medical management. Prognostic factors include the patient's overall health status, nutritional status, the specific surgical technique employed, and the presence of comorbidities. With proper surgical technique, comprehensive preoperative and postoperative care, and dedicated ostomy nursing support, most patients with an ileostomy can lead full, active, and productive lives, often experiencing a significant improvement in their health status compared to their pre-surgical condition.

\section*{References}
\begin{enumerate}
  \item MedlinePlus. Ileostomy. U.S. National Library of Medicine; 2024. Available from: https://medlineplus.gov/ileostomy.html
  
  \item Brunicardi F, Andersen DK, Billiar TR, et al. Schwartz's Principles of Surgery. 11th ed. McGraw-Hill Education; 2019.
  
  \item Fazio VW, Church JM, Fleshman JW, et al. Current Therapy in Colon and Rectal Surgery. 3rd ed. Elsevier; 2017.
  
  \item American Society of Colon and Rectal Surgeons (ASCRS). Clinical Practice Guidelines for Ostomy Surgery. 2023.
\end{enumerate}

\clearpage
